%! TeX program = lualatex
\documentclass[11pt,oneside,openany]{memoir}
\input{/Users/gcrescenzo/Documents/School/LaTeX Documents/latex-class-notes/preamble-notes.tex}
\usepackage{kbordermatrix}
\usepackage{nicematrix}
\input{/Users/gcrescenzo/Documents/School/LaTeX Documents/latex-class-notes/makros.tex}
\renewcommand*{\emph}[1]{\textit{#1}}

\renewcommand{\projecttitle}{Linear Algebra \& Module Theory}

\begin{document}

\pagenumbering{roman}
\chapter*{Preface}
\vspace{-35pt}
\noindent After introducing categories, we will study modules, which are vector spaces over commutative rings. These are the basis for a proper treatment of linear algebra, but also for the study of sheaf cohomology in algebraic geometry. We will develop the deep and great idea of universal properties, which changed the way we think about mathematics. The main theorem of the course is the fundamental theorem of finitely generated modules over a PID, which classifies finitely generated abelian groups and provides the correct theoretical framework for computing canonical forms of linear transformations. We will spend the rest of the course on applications of these ideas to bilinear forms and matrix Lie groups. Snarf

\vspace{-35pt}
\tableofcontents
\vfill
\specialdate
Last update: \today

\pagenumbering{arabic}

\chapter{Modules Over a Commutative Ring}

\section{Categories}

\subsection{Basic Definitions}

\begin{definition}
	A \emph{category} $\scrC$ is a pair of classes $(\obj(\scrC), \Hom(\scrC))$, where $\obj(\scrC)$ are the \emph{objects} of $\scrC$, and $\Hom(\scrC)$ are the \emph{morphisms} of $\scrC$. Every morphism involes two objects, the \emph{source} (i.e., the domain of the morphism) and the \emph{target} (i.e., the codomain of the morphism). Every object must be identified with a unique \emph{identity morphism} sending the object to iteslf. The morphisms must satisfy composition, and such a composition must be associative. The set of morphisms between two objects $A$ and $B$ is denoted $\Hom_\scrC(A,B)$.
\end{definition}

\begin{example}
	The category \catname{Set} has as objects sets and morphisms set maps. The category \catname{Grp} and \catname{Ring} are defined similarly. The category \catname{Top} has as objects topological spaces and morphisms continuous maps. The category \catname{$k$-Vec} has as objects $k$-vector spaces and morphisms linear transformations.
\end{example}

\begin{definition}
	Let $\scrC$ be a category.
	\mbox{}
	\begin{enumerate}[label = (\arabic*),itemsep=1pt,topsep=3pt]
		\item An \emph{initial object} of $\scrC$ is some $I \in \obj(\scrC)$ such that for every $A \in \obj(\scrC)$, there exists a unique morphism $I \rightarrow A$.
		\item A \emph{terminal object} of $\scrC$ is some $T \in \obj(\scrC)$ such that for every $B \in \obj(\scrC)$, there exists a unique morphism $B \rightarrow T$.
		\item A \emph{zero object} of $\scrC$ is an object that is both initial and terminal.
	\end{enumerate}
\end{definition}

\begin{example}
	In \catname{Set}, the emptyset is an initial object and the singletons are terminal objects. There does not exist a zero object in \catname{Set}. In \catname{Ring}, \catname{Grp}, \catname{$k$-Vec}, $\{0\}$ is a zero object.
\end{example}

\subsection{Products and Coproducts}


\begin{definition}
	Let $\scrC$ be a category and $\{A_i\}_{i \in I}$ a family of objects in $\scrC$.
	\mbox{}
	\begin{enumerate}[label = (\arabic*),itemsep=1pt,topsep=3pt]
		\item A \emph{cone} over $\{A_i\}_{i \in I}$ is a pair $\bigl( A,\{f_i\}_{i \in I} \bigr)$ such that $f_i:A \rightarrow A_i$ is a morphism in $\scrC$ for all $i$.
		\item A \emph{cocone} over $\{A_i\}_{i \in I}$ is a pair $\bigl( B, \{g_i\}_{i \in I} \bigr)$ such that $g_i:A_i \rightarrow B$ is a morphism in $\scrC$ for all $i$.
	\end{enumerate}
\end{definition}

\begin{example}
	Let $\scrC$ be a category and let $\{A_i\}_{i \in I}$ be a family of objects in $\scrC$. The category \catname{Cone$(\{A_i\}_{i \in I})$} (or just denoted \catname{Cone} if the context is clear) has as objects cones over $\{A_i\}_{i \in I}$ and as morphisms $\bigl( A,\{f_i\}_{i \in I} \bigr) \rightarrow \bigl( A', \{g_i\}_{i \in I} \bigr)$ making the following diagram commute for each $i$:
	\begin{center}
		\begin{tikzcd}
A \arrow[rr] \arrow[rd, "f_i"'] &     & A' \arrow[ld, "g_i"] \\
                                & A_i &                     
\end{tikzcd}
	\end{center}
\end{example}

\begin{example}
	Let $\scrC$ be a category and let $\{B_i\}_{i \in I}$ be a family of objects in $\scrC$. The category \catname{Cocone$(\{B_i\}_{i \in I})$} (or just denoted \catname{Cocone} if the context is clear) has as objects cocones over $\{B_i\}_{i \in I}$ and as morphisms $\bigl( B,\{f_i\}_{i \in I} \bigr) \rightarrow \bigl( B', \{g_i\}_{i \in I} \bigr)$ making the following diagram commute for each $i$:
	\begin{center}
		\begin{tikzcd}
             & B_i \arrow[ld, "f_i"'] \arrow[rd, "g_i"] &    \\
B \arrow[rr] &                                          & B'
\end{tikzcd}
	\end{center}
\end{example}

\begin{definition}
	Let $\scrC$ be a category and let $\{A_i\}_{i \in I}$ be a family of objects in $\scrC$. The \emph{product} over $\{A_i\}_{i \in I}$, denoted by $\bigl( \prod_{i \in I}A_i, \{\pi_i\}_{i \in I} \bigr)$, is the terminal object in \catname{Cone$(\{A_i\}_{i \in I})$}. 
\end{definition}

\begin{proposition}
	Let $\scrC$ be a category and let $\{A_i\}_{i \in I}$ be a family of objects in $\scrC$. If the product of $\{A_i\}_{i \in I}$ exists, then for any cone $\bigl( A, \{\tau_i\}_{i \in I} \bigr) \in \obj(\catname{Cone})$, there is a unique map $\varphi:A \rightarrow \prod_{i \in I}A_i$ making the following diagram commute for all $i$:
	\begin{center}
		\begin{tikzcd}
                                                              & \prod_{i \in I}A_i \arrow[d, "\pi_i"] \\
A \arrow[r, "\tau_i"'] \arrow[ru, "\exists! \varphi", dotted] & A_i                                  
\end{tikzcd}
	\end{center}
\end{proposition}

\begin{definition}
	Let $\scrC$ be a category and let $\{B_i\}_{i \in I}$ be a family of onjects in $\scrC$. The \emph{coproduct} over $\{B_i\}_{i \in I}$, denoted by $\bigl( \coprod_{i \in I}B_i, \{\iota_i\}_{i \in I} \bigr)$, is the initial object in \catname{Cocone$(\{B_i\}_{i \in I})$}.
\end{definition}

\begin{proposition}\label{prop:cocone-up}
	Let $\scrC$ be a category and let $\{B_i\}_{i \in I}$ be a family of objets in $\scrC$. If the coproduct of $\{B_i\}_{i \in I}$ exists, then for any cocone $\bigl( B, \{\theta_i\}_{i \in I} \bigr) \in \obj(\catname{Cocone})$, there is a unique map $\varphi:\coprod_{i \in I} B_i \rightarrow B$ making the following diagram commute for all $i$:
	\begin{center}
		\begin{tikzcd}
\coprod_{i \in I}B_i \arrow[rd, "\exists!\varphi", dotted] &   \\
B_i \arrow[r, "\theta_i"] \arrow[u, "\iota_i"]             & B
\end{tikzcd}
	\end{center}
\end{proposition}

\begin{example}
	Consider the collection $\{\bbZ_n\}_{n \in \bbN} \subseteq \obj(\catname{Grp})$. Let $\pi_n :\bbZ \rightarrow \bbZ_n$ be the natural quotient maps. There does not exist a map of cones $\bigl( \bbZ , \{\pi_n\}_{n \in \bbN}\bigr)\rightarrow 0$. Indeed, the following diagram does not commute:
	\begin{center}
		\begin{tikzcd}
\bbZ \arrow[rd, "\pi_n"'] \arrow[rr] &        & 0 \arrow[ld, "0"] \\
                                     & \bbZ_n &                  
\end{tikzcd}
	\end{center}
	However, there does exist a map of cones $0 \rightarrow \bigl( \bbZ, \{\pi_n\}_{n \in \bbN} \bigr)$. Indeed, the diagram below commutes:
	\begin{center}
		\begin{tikzcd}
0 \arrow[rd, "0"'] \arrow[rr] &        & \bbZ \arrow[ld, "\pi_n"] \\
                              & \bbZ_n &                         
\end{tikzcd}
	\end{center}
\end{example}

\subsection{Exact Sequences}

\begin{definition}
	Let $\scrC$ be an abelian category (e.g., \catname{Grp}, \catname{Ring}, etc). A sequence:
	\begin{center}
		\begin{tikzcd}
A_0 \arrow[r, "f_1"] & A_1 \arrow[r, "f_2"] & A_2 \arrow[r, "f_3"] & ... \arrow[r, "f_n"] & A_n
\end{tikzcd}
	\end{center}
	of objects and morphisms in $\scrC$ is \emph{exact at $A_i$} if $\Image(f_i) = \ker(f_{i + 1})$. The sequence is called \emph{exact} if it is exact at each $A_i$. A \emph{short exact sequence} is an exact sequence of the form:
	\begin{center}
		\begin{tikzcd}
0 \arrow[r] & A \arrow[r, "f"] & B \arrow[r, "g"] & C \arrow[r] & 0
\end{tikzcd}
	\end{center}
	for objects $A,B,C$ in $\scrC$ and morphisms $f,g$.
\end{definition}

\begin{example}
	Let $G,H \in \obj(\catname{Grp})$ and suppose $H$ is a subgroup of $G$. Let $\iota:H \rightarrow G$ be the inclusion map and $\pi: G \rightarrow G/H$ the quotient projection. Then the following sequence is exact:
	\begin{center}
		\begin{tikzcd}
0 \arrow[r] & H \arrow[r, "\iota"] & G \arrow[r, "\pi"] & G/H \arrow[r] & 0
\end{tikzcd}
	\end{center}
\end{example}

\begin{lemma}\label{lemma:splitting-lemma}
	Let $\scrC$ be an abelian category and:
	\begin{center}
		\begin{tikzcd}
0 \arrow[r] & A \arrow[r, "f"] & B \arrow[r, "g"] & C \arrow[r] & 0
\end{tikzcd}
	\end{center}
	a short exact sequence of objects in $\scrC$. The following are equivalent:
	\mbox{}
	\begin{enumerate}[label = (\arabic*),itemsep=1pt,topsep=3pt]
		\item There exists a morphism $\rho:B \rightarrow A$ such that $\rho \circ f = \id_A$.
		\item There exists a morphism $\sigma:C \rightarrow B$ such that $g \circ \sigma = \id_C$.
		\item There is an isomorphism $\varphi:B \rightarrow A \oplus C$ such that $\varphi \circ f = \iota_A$ and $g \circ \varphi^{-1} = \pi_C$, where $\iota_A$ is the natural inclusion of $A$ into $A \oplus C$ and $\pi_C$ is the natural projection of $A \oplus C$ into $C$. In particular, the following diagram commutes:
			\begin{center}
\begin{tikzcd}
            &                                          & B \arrow[dd, "\varphi", dotted] \arrow[rd, "g"] &             &   \\
0 \arrow[r] & A \arrow[ru, "f"] \arrow[rd, "\iota_A"'] &                                                 & C \arrow[r] & 0 \\
            &                                          & A \oplus C \arrow[ru, "\pi_C"']                 &             &  
\end{tikzcd}
			\end{center}
	\end{enumerate}
\end{lemma}
	\begin{proof}
		(1)$\Rightarrow$(3) Note that $\ker \rho + \Image f \subseteq B$ automatically. For any $b \in B$, we have:
		\[
			b = \bigl( b - f(\rho(b)) \bigr) + f(\rho(b)).
		\] 
		Since $b - f(\rho(b)) \in \ker \rho$ and $f(\rho(b)) \in \Image f$, this gives $\ker \rho + \Image f = B$. Now let $x \in \ker \rho \cap \Image f$. Then $\rho(x) = 0$ and $x = f(a)$ for some $a \in A$. But this means $0 = \rho(x) = \rho(f(a)) = a$. So $x = f(0) = 0$. Thus $B = \ker \rho \oplus \Image f$. We will now show that $\restr{g}{\ker \rho}:\ker \rho \rightarrow C$ is an isomorphism. By exactness, note that $g$ is surjective. Hence for any $c \in C$, there exists $b \in B$ with $g(b) = c$. By the above decomposition, we have:
		\[
			g(b - f(\rho(b)) = g(b) - g(f(\rho(b))) = c.
		\] 
		Thus $\restr{g}{\ker \rho}$ is surjective. Now if $k \in \ker \rho$ and $g(k) = 0$, then again by exactness $k \in \ker g = \Image f$. Thus $k \in \ker \rho \cap \Image f = \{0\}$, giving $\ker \rho \simeq C$. We also have that $f:A \rightarrow \Image f$ is an isomorphism. Indeed, any map always surjects onto its image, and $\rho \circ f = \id_A$ makes it injective. Thus $B \simeq A \oplus C$. The particular isomorphism $\varphi:B \rightarrow A \oplus C$ is given by $\varphi(b) := (\rho(b),g(b))$. We finally have $\varphi(f(a)) = (a,0) = \iota_A(a)$, and $\pi_C(\varphi(b)) = g(b)$, establishing (3).
	\end{proof}

\begin{definition}
	A short exact sequence satisfying Lemma~\ref{lemma:splitting-lemma} is called a \emph{split exact sequence}.
\end{definition}

\section{Modules}

Unless otherwise stated, assume $R$ to be a commutative ring with $1$.

\subsection{Definitions and Examples}

\begin{definition}
	A \emph{left $R$-module} is an abelian group $(M,+)$ together with an action:
	\[
		R \times M \rightarrow M, \hspace{2pt} (r,m) \mapsto rm
	\] 
	which satisfies, for all $r,s \in R$ and $m \in M$:
	\mbox{}
	\begin{enumerate}[label = (\arabic*),itemsep=1pt,topsep=3pt]
		\item $(r+s)m = rm + sm$;
		\item $r(sm) = (rs)m$;
		\item $r(m+n) = rm + rn$.
	\end{enumerate}
	A \emph{right $R$-module} is defined similarly. If $M$ is both a left and right $R$-module, then we simply refer to it as an $R$-module. If $M$ is instead a ring $(M,+,\cdot)$ and the action $R \times M \rightarrow M$ also satisfies, for all $r \in R$ and $m,n \in M$:
	\mbox{}
	\begin{enumerate}[label = (\arabic*),itemsep=1pt,topsep=3pt]
		\addtocounter{enumi}{3}
		\item $r(mn) = (rm)n = m(rn)$,
	\end{enumerate}
	then we say $M$ is an \emph{$R$-algebra} (resp., \emph{left} or \emph{right}).
\end{definition}

\begin{example}
	Every ring is a module over itself. If $I \subseteq R$ is an ideal, then $I$ and $R/I$ are both $R$-modules.
\end{example}

\begin{definition}
	Let $M$ and $N$ be $R$-modules. An \emph{$R$-linear homomorphism} is a group homomorphism $f:M \rightarrow N$ which also satisfies $f(rm) = r f(m)$ for all $r \in R$. If an $R$-linear homomorphism $M \rightarrow N$ is bijective, then we call it an \emph{isomorphism} and write $M \simeq N$. The set of all $R$-linear homomorphisms from $M$ to $N$ is denoted $\Hom_R(M,N)$. The set of all $R$-isomorphisms from $M$ to $N$ is denoted $\Aut_R(M,N)$. If $M = N$, we denote $\Hom_R(M,N)$ by $\End_R(M)$ and $\Aut_R(M,N)$ by $\Aut_R(M)$.
\end{definition}

\begin{example}
	Let $M$ and $N$ be $R$-modules. Then $\Hom_R(M,N)$ is an $R$-module, $\End_R(M)$ is an $R$-algebra, and $\Aut_R(M)$ is a group.
\end{example}

\iffalse
\begin{example}
	Consider $\bbZ_m$ and $\bbZ_n$ as $\bbZ$-modules. We wish to compute $\Hom_{\bbZ}(\bbZ_m,\bbZ_n)$. Note that each $\varphi \in \Hom_{\bbZ}(\bbZ_m,\bbZ_n)$ is completely determined by $\varphi(1)$ because $\varphi(n) = n \varpi(1)$. {\color{red} DNF}
\end{example}
\fi

\begin{definition}
	$R$-modules and $R$-module homomorphisms form a category, denoted by \emph{\catname{$R$-Mod}}.
\end{definition}

\begin{example}
	If $R$ is a field, then $\catname{$R$-Mod} = \catname{$R$-Vec}$.
\end{example}

\begin{example}
	Let $A \in \obj(\catname{AbGrp})$. For any $a \in A$, recall that:
	\[
		na = \begin{cases} a + \ldots + a \text{ ($n$ times) }, & n > 0 \\ 0, & n =0 \\ -a - \ldots - a \text{ ($-n$ times) }, & n < 0\end{cases}.
	\]
	This defines an action of $\bbZ$ on $A$, making $A$ into a $\bbZ$-module. Hence $\catname{AbGrp} \subseteq \catname{$\bbZ$-Mod}$. Conversely, if $M$ is any $\bbZ$-module then $M$ is automatically an abelian group. Hence $\catname{AbGrp} = \catname{$\bbZ$-Mod}$.
\end{example}

\begin{example}
	Let $k$ be a field and consider the polynomial ring $k[x]$. Let $V$ be a $k$-vector space {\color{red} DNF}
\end{example}

\begin{definition}
	Given an $R$-module $M$, the \emph{annihilator} of $M$ is $\Ann_R(M) := \{a \in R : am= 0 \text{ for all $m \in M$ }\}$.
\end{definition}

\begin{proposition}
	$\Ann_R(M) \subset R$ is a two-sided ideal, and we have a commutative diagram:
	\begin{center}
		\begin{tikzcd}
R \times M \arrow[d] \arrow[r] & M \\
R/\Ann_R(M) \times M \arrow[ru]         &  
\end{tikzcd}
	\end{center}
\end{proposition}
	\begin{proof}
	For convenience let $A:= \Ann_R(M)$. Note that $A \neq \emptyset$ since $0 \in A$. Let $x,y \in A$. Then $xm = 0$ and $ym = 0$ for all $m \in M$. Hence $(x-y)m = xm - ym = 0$, so $A$ is a subgroup of $R$. Now if $x \in A$ and $r \in R$, we have $(xr)m = (rx)m = r(xm) = 0$. Similarly, $(rx)m = r(xm) = 0$. Thus $xr,rx \in A$, hence $A$ is a two-sided ideal.

	Let $R \times M \rightarrow M$ be the usual scalar multiplication, and define $R \times M \rightarrow R/A \times M$ by $(r,m) \mapsto (r + A,m)$. Similarly, define $R/A \times M \rightarrow M$ by $(r + A,m) \mapsto rm$. We must show that this map is well-defined. Suppose $(r_1 + a,m_1) = (r_2 + a,m_2)$. Then $m_1 = m_2$ and $r_1 + A = r_2 + A$. Hence $r_1-r_2 \in A$, so $(r_1-r_2)m = 0$ for all $m \in M$. Thus $(r_1-r_2)m_1 = r_1m_1 - r_2m_1 = r_1m_1 - r_2m_2 = 0$. Thus $r_1m_1 = r_2m_2$, so the map is well-defined. With this, it is clear that these maps make the above diagram commute. 
\end{proof}

\begin{example}
	Let $V := \bbZ_{24} \times \bbZ_{15} \times \bbZ_{50}$. We will compute $\Ann_{\bbZ}(V) = \{z \in \bbZ : z(a,b,c) = 0 \text{ for all $(a,b,c) \in V$}\}$. Intuitively, if $z \in \bbZ$ is the least common multiple of $24$, $15$, and $50$ then it should kill any coset of $V$. Claim: $\Ann_{\bbZ}(V) = \left<\lcm(24,15,50) \right>$. If $z \in \left<\lcm(24,15,50) \right>$, then as was just explained $z \in \Ann_{\bbZ}(V)$, hence $\left<\lcm(24,15,50) \right> \subseteq \Ann_{\bbZ}(V)$. If $z \in \Ann_{\bbZ}(V)$, then $z(a,b,c) = 0$ for all $(a,b,c) \in V$. But $(a,b,c) = a(1,0,0) + b(0,1,0) + c(0,0,1)$. Hence it is both necessary and sufficient for $z(1,0,0) = 0$, $z(0,1,0) = 0$, and $z(0,0,1)$. Equivalently, $z \equiv 0 \pmod{24}$, $z \equiv 0 \pmod{15}$, and $z \equiv 0 \pmod{50}$. This is equivalent to $z \equiv 0 \pmod{\lcm(24,15,50)}$, thus $\Ann_{\bbZ}(V) = \left<\lcm(24,15,50) \right>$.
\end{example}

\begin{definition}
	Let $M$ be an $R$-module.
	\mbox{}
	\begin{enumerate}[label = (\arabic*),itemsep=1pt,topsep=3pt]
		\item A \emph{submodule} $N \subseteq M$ is a subgroup that is itself an $R$-module.
		\item $M$ is \emph{generated} by a subset $X \subset M$ if $M = \{\sum_{i=1}^{n} a_i x_i : a_i \in R, x_i \in X, n \in \bbN\}$. In this case, we write $M = \left<X \right>$.
		\item $M$ is \emph{finitely generated} if $M = \left<X \right>$ for a finite set $X \subset M$.
		\item $M$ is \emph{cyclic} if $M = \left<m \right>$ for some $m \in M$.
		\item $M$ is \emph{irreducible} (or \emph{simple}) if it has no submodules.
	\end{enumerate}
\end{definition}

\begin{proposition}
	If $M$ is an irreducible $R$-module, then $\Ann_R(M)$ is a maximal ideal of $R$.
\end{proposition}
	\begin{proof}
		Suppose $b \not\in \Ann_R(M)$. Then $bx \neq 0$ for some $x \in M$. Hence $\left<bx \right>$ is a submodule of $M$. But since $M$ is irreducible and $bx$ is nonzero, it must be the case that $\left<bx \right> = M$. But since $x \in M$, we can write $x = cbx$ for some $c \in R$. Thus $0 = (cb-1)x$, hence $cb - 1 \in \Ann_R(M)$. Equivalently, $cb + \Ann_R(M) = 1 + \Ann_R(M)$, which means $b$ admits an inverse, namely $c + \Ann_R(M)$. Since every element of $R/\Ann_R(M)$ is invertible, it is a field. Thus $\Ann_R(M)$ is a maximal ideal of $R$. 
	\end{proof}

\subsection{Isomorphism Theorems}

{\color{red} DNF}

\subsection{Free Modules \& Linear Algebra}

\begin{definition}
	A \emph{basis} for an $R$-module $M$ is a set $\{x_i\}_{i \in I} \subset M$ such that $\{x_i\})_{i \in I}$ is linearly independent and spans $M$. In this case, we say $M$ is \emph{free}. 
\end{definition}

\begin{theorem}\label{thm:coprod-exist}
	Coproducts exist in \catname{$R$-Mod}.
\end{theorem}
	{\color{red}
	\begin{proof}
			
	\end{proof}}

\begin{definition}
	Let $I$ be an index set. The \emph{free module} on $I$ is $R^{(I)} := \coprod_{I}R_i$. If $1_i$ denotes the mutliplicative identity of $R_i$, then $R^{(I)}$ has $\{1_i\}_{i \in I}$ as a basis. An $R$-module $M$ is called \emph{free} if there exists an index set $I$ such that $M \simeq R^{(I)}$.
\end{definition}

\begin{corollary}
	An $R$-module $M$ is free if and only if there exists an index set $I$ such that $M \simeq R^{(I)}$.
\end{corollary}

\begin{example}
	If $k$ is a field, then any $k$-vector space is a free $k$-module. For every $n \in \bbN$, the set $R^n$ is a free $R$-module. The set $R[x]$ is also a free module, since it has $\{x^n\}_{n \in \bbN}$ as a basis.
\end{example}

\begin{theorem}
	Let $I$ be an index set, and $M$ and $R$-module. Suppose $f:I \rightarrow M$ is any set map and $\iota : I \rightarrow R^{(I)}$ is the natural inclusion sending $i \mapsto 1_i$. Then there exists a unique $R$-linear homomorphism $g:R^{(I)} \rightarrow M$ making the following diagram commute:
	\begin{center}
		\begin{tikzcd}
I \arrow[rd, "f"'] \arrow[r, "\iota"] & R^{(I)} \arrow[d, "\exists!g", dotted] \\
                                      & M                                     
\end{tikzcd}
	\end{center}
\end{theorem}
	\begin{proof}
		This follows immediately from Theorem~\ref{thm:coprod-exist} Proposition~\ref{prop:cocone-up}.
	\end{proof}

\begin{proposition}\label{prop:hom-isom-mat}
	There is an isomorphism $\Hom_R(R^n,R^m) \simeq \Mat_{m \times n}(R)$.
\end{proposition}
	\begin{proof}
		If $\{1_i\}_{i = 1}^n$ is a basis for $R^n$, define $\Hom_R(R^n,R^m) \rightarrow \Mat_{m \times n}(R)$ by: $$\varphi \mapsto \begin{bmatrix} \varphi(1_1) & \varphi(1_2) & \ldots &\varphi(1_n) \end{bmatrix}.$$ Properties of vectors and the linearity of $\varphi$ gives that this mapping is linear. This map is also clearly injective: if two maps $\varphi_1$ and $\varphi_2$ agree on all basis elements, then they must be equal. Now if $A \in \Mat_{m \times n}(R)$, let $\varphi_A:R^n \rightarrow R^m$ be given by $\varphi_A(x) \mapsto Ax$. Then:
		\[
			\varphi_A \mapsto \begin{bmatrix} A 1_1 & A 1_2 & \ldots & A 1_n \end{bmatrix} = A.
		\] 
		Thus $\Hom_R(R^n,R^m)\simeq \Mat_{m \times n}(R)$.
	\end{proof}

\begin{definition}
	\mbox{}
	\begin{enumerate}[label = (\arabic*),itemsep=1pt,topsep=3pt]
		\item The \emph{determinant} is a map $\det:\Mat_n(R) \rightarrow R$ given by $(a_{i,j}) \mapsto \sum_{\sigma \in S_n}^{}(-1)^\sigma a_{1,\sigma(1)}a_{2,\sigma(2)}\ldots a_{n,\sigma(n)}$.
		\item The \emph{cofactor} of $(a_{ij})$ is $(-1)^{i + j}\Det (m_{ij})$, where $(m_{ij})$ is the submatrix of $(a_{ij})$ with row $i$ and column $j$ deleted.
		\item The \emph{adjugate} (or \emph{adjoint}) of $(a_{ij})\in \Mat_n(R)$ is $\adj\bigl( (a_{ij}) \bigr) := (b_{ij})$, where $b_{ij} = a_{ji}$.
	\end{enumerate}
\end{definition}

\begin{proposition}
	Let $A := (a_{ij})$ and $B := (b_{ij})$ be elements of $\Mat_n(R)$.
	\mbox{}
	\begin{enumerate}[label = (\arabic*),itemsep=1pt,topsep=3pt]
		\item $\det AB = \det A \det B$.
		\item $(\det A)I_n = A \adj(A) = \adj(A)A$.
		\item $\det AB = \det BA$.
		\item $\det$ is $R$-multilinear on rows and columns.
		\item $\det$ is alternating on rows and columns.
		\item $\det I_n = 1$.
	\end{enumerate}
\end{proposition}
	{\color{red}
	\begin{proof}
			
	\end{proof}}

\begin{proposition}
	There is an isomorphism $\Aut_R(R^n) \simeq \GL_n(R) = \{A \in \Mat_n(R) : \det A \in R^\times\}$. Moreover, $\det : GL_n(R) \rightarrow R^\times$ is a group homomorphism.
\end{proposition}
	\begin{proof}
		If $\{1_i\}_{i=1}^n$ is a basis for $R^n$, define $\Phi: \Aut_R(R^n) \rightarrow \Mat_n(R)$ by:
		\[
			\varphi \mapsto \begin{bmatrix} \varphi(1_1) & \varphi(1_2) & \ldots & \varphi(1_n) \end{bmatrix}.
		\]
		This map is injective by Proposition~\ref{prop:hom-isom-mat}. For any $\Phi(\varphi) \in \Image(\Phi)$, note that:
		\begin{align*}
			1
			&= \det I_n \\
			&= \det \Phi(\varphi \varphi^{-1}) \\
			&= \det \bigl(\Phi(\varphi) \Phi(\varphi)^{-1}\bigr) \\
			&= \det \Phi(\varphi) \det \Phi(\varphi)^{-1} \\
		\end{align*}
		Hence $\Image(\Phi) \subseteq \GL_n(R)$. Conversely, if $A \in \GL_n(R)$, then again by Proposition~\ref{prop:hom-isom-mat} we have $A \in \Image(\Phi)$. Hence $\Image(\Phi) = \GL_n(R)$, so by the First Isomorphism Theorem $\Aut_R(R^n) \simeq \GL_n(R)$.
	\end{proof}

\begin{theorem}[Cayley-Hamilton]
	Let $M$ be a finitely generated $R$-module. Then each $\varphi \in \End_R(M)$ is a root of a monic in $R[x]$. 
\end{theorem}
	\begin{proof}
		Suppose $M = \left<e_1,\ldots ,e_n \right>$. Since $\varphi(e_i) \in M$, there exists $a_{1i},a_{2i},\ldots ,a_{ni}$ such that $\varphi(e_i) = \sum_{\alpha =1}^{n} a_{\alpha i}e_{\alpha}$. So for each $i = 1,\ldots ,n$ we have:
		\begin{align*}
\varphi(e_1) &= \sum a_{\alpha 1}e_\alpha,\\
\varphi(e_2) &= \sum a_{\alpha 2}e_\alpha,\\
&\vdots\\
\varphi(e_n) &= \sum a_{\alpha n}e_\alpha.
\end{align*}

We can rearrange each equation as follows:
\begin{align*}
(\varphi-a_{11})(e_1)-a_{21}e_1-\cdots-a_{n1}e_1 &= 0\\
a_{12}e_2-(\varphi-a_{22})(e_2)-a_{32}e_2-\cdots-a_{n2}e_2 &= 0\\
&\vdots\\
a_{1n}e_n-a_{2n}e_n-\cdots-a_{(n-1)n}e_n-(\varphi-a_{nn})(e_n) &= 0.
\end{align*}

We can express the above system this using a matrix as follows:
\begin{align*}
\begin{bmatrix}
\varphi-a_{11} & -a_{12} & -a_{13} & \cdots & -a_{1n}\\
-a_{21} & \varphi-a_{22} & & & \\
-a_{31} & & \ddots & & \\
\vdots & & & & \\
-a_{n1} & & & & \varphi-a_{nn}
\end{bmatrix}
\begin{bmatrix}
	e_1\\\phantom{a} \\
	\vdots\\\phantom{a}\\
e_n
\end{bmatrix}
&= 0.
\end{align*}
Let $A_\varphi:R^n \rightarrow R^n$ denote the above matrix, and formally $A_\varphi \in \Mat_n(R[\varphi])$. By construction, $A_\varphi(1_i e_i) = 0$ for each $i=1,\ldots ,n$. Hence $(\adj A_\varphi)A_\varphi(1_i e_i) = 0$ for each $i=1,\ldots ,n$. Thus $(\adj A_\varphi \cdot A_\varphi)(1_i e_i) = (\det A_\varphi)(1_i e_i) = 0$ for each $i=1,\ldots ,n$. It must be the case that $\det A_\varphi = 0$. Thus $\varphi$ is a solution to the monic polynomial $\det A_x \in R[x]$.
	\end{proof}

\begin{example}
	Consider the $\bbZ$-algebra $\bbZ[i]$. This module is finitely generated by $1$ and $i$. For each $a \in \bbZ[i]$, define $L_a :\bbZ[i] \rightarrow \bbZ[i]$ by $z \mapsto az$. By the Cayley-Hamilton Theorem, there exists a monic $f \in \bbZ[x]$ such that $f(L_a) = 0$. As an example, consider $L_{3+2i}$, we have:
	\begin{align*}
T_{3+2i}(1) &= 3+2i,\\
T_{3+2i}(i) &= -2+3i.
\end{align*}

Hence:
\begin{align*}
[T] &= \begin{bmatrix}
3 & -2\\
2 & 3
\end{bmatrix}.
\end{align*}

Following the construction in the proof of the Cayley--Hamilton Theorem, let:
\begin{align*}
A_x &:= \begin{bmatrix}
3-x & -2\\
2 & 3-x
\end{bmatrix}.
\end{align*}

Then $\det A_x=x^2-6x+13$, and hence $L_{3+2i}^{2}-6L_{3+2i}+13=0$.
\end{example}


\iffalse
\begin{appendices}
\include{app2.tex}
\include{app1.tex}
\end{appendices}
\fi

\end{document}

